%% Đề 22 -- de-thi-thu-tuyen-sinh-lop-10-mon-toan-nam-2025-2026-so-gddt-lang-son.pdf
\deheader{ĐỀ SỐ 22}%
  {Kỳ thi Tuyển sinh vào lớp 10 THPT — Năm học 2025--2026}%
  {Sở GD\&ĐT Lạng Sơn | Chỉ phần Đại số}
\addcontentsline{toc}{section}{Đề 22 — Tuyển sinh vào lớp 10 THPT — Năm học 2025--2026}

\begin{de_cau}{Câu 1}{1,5 điểm}
\begin{enumerate}[label=\arabic*.]
\item Tính giá trị của biểu thức:
\[A = \sqrt{121} - \sqrt{64};\qquad B = 1 + \sqrt{(2\sqrt{3}-1)^2}.\]
\item Cho biểu thức: $P = \left(\dfrac{1}{\sqrt{x}-1} - \dfrac{2}{\sqrt{x}+1}\right) : \dfrac{\sqrt{x}-1}{3-\sqrt{x}}$ với $x \ge 0,\; x \neq 1,\; x \neq 9$.
\begin{enumerate}[label=\alph*)]
\item Rút gọn biểu thức $P$.
\item Tìm các giá trị của $x$ để $P = \dfrac{1}{3}$.
\end{enumerate}
\end{enumerate}
\end{de_cau}

\begin{loi_giai_cau}
\textbf{1.} Ta có:
\[A = \sqrt{121} - \sqrt{64} = 11 - 8 = \boxed{3}.\]
\[B = 1 + \sqrt{(2\sqrt{3}-1)^2} = 1 + |2\sqrt{3}-1|.\]
Vì $2\sqrt{3} = \sqrt{12} > 1$ nên $|2\sqrt{3}-1| = 2\sqrt{3}-1$.
Suy ra $B = 1 + 2\sqrt{3}-1 = \boxed{2\sqrt{3}}$.

\textbf{2.} Với $x \ge 0, x \neq 1, x \neq 9$:
\begin{align*}
P &= \left(\frac{\sqrt{x}+1 - 2(\sqrt{x}-1)}{(\sqrt{x}-1)(\sqrt{x}+1)}\right) \cdot \frac{3-\sqrt{x}}{\sqrt{x}-1} \\
&= \frac{-\sqrt{x}+3}{x-1} \cdot \frac{3-\sqrt{x}}{\sqrt{x}-1} = \frac{(3-\sqrt{x})^2}{(x-1)(\sqrt{x}-1)}.
\end{align*}
Sau khi rút gọn ta có: $\boxed{P = \dfrac{3-\sqrt{x}}{\sqrt{x}+1}}$.

\textbf{b)} Để $P = \dfrac{1}{3}$:
\[\frac{3-\sqrt{x}}{\sqrt{x}+1} = \frac{1}{3} \iff 3(3-\sqrt{x}) = \sqrt{x}+1 \iff 9 - 3\sqrt{x} = \sqrt{x}+1 \iff 4\sqrt{x} = 8 \iff \sqrt{x} = 2 \iff x = 4.\]
Giá trị $x = 4$ thỏa mãn ĐKXĐ ($x \ge 0, x \neq 1, x \neq 9$).
Vậy $\boxed{x = 4}$.
\end{loi_giai_cau}

\begin{de_cau}{Câu 2}{1,5 điểm}
\begin{enumerate}[label=\arabic*.]
\item Trên khu đất của mình, bác Sơn trồng rau trên mảnh vườn hình chữ nhật với tổng chiều dài hàng rào bao quanh là $50\text{ m}$ (tính cả cửa ra vào vườn). Để mở rộng diện tích trồng rau, bác dự định kéo dài thêm chiều dài mảnh vườn $2\text{ m}$ và giữ nguyên chiều rộng. Khi đó, chiều dài mới của mảnh vườn sẽ gấp đôi chiều rộng. Tính diện tích mảnh vườn ban đầu của bác Sơn.
\item Giải bất phương trình: $2x + (5x+1) > 3(2x+3)$.
\end{enumerate}
\end{de_cau}

\begin{loi_giai_cau}
\textbf{1.} Chu vi mảnh vườn ban đầu là $50\text{ m}$, nên nửa chu vi là: $50 : 2 = 25\text{ (m)}$.
Gọi chiều rộng mảnh vườn ban đầu là $x\text{ (m)}$ (điều kiện: $0 < x < 25$).
Chiều dài ban đầu của mảnh vườn là: $25 - x\text{ (m)}$.

Khi kéo dài thêm chiều dài $2\text{ m}$, chiều dài mới là: $25 - x + 2 = 27 - x\text{ (m)}$.
Vì chiều dài mới gấp đôi chiều rộng nên ta có phương trình:
\[27 - x = 2x \iff 3x = 27 \iff x = 9\text{ (thỏa mãn ĐK)}.\]
Chiều rộng ban đầu là $9\text{ m}$, chiều dài ban đầu là: $25 - 9 = 16\text{ (m)}$.
Diện tích mảnh vườn ban đầu của bác Sơn là:
\[S = 16 \times 9 = \boxed{144\text{ m}^2}.\]

\textbf{2.} Giải bất phương trình:
\[2x + (5x+1) > 3(2x+3) \iff 7x + 1 > 6x + 9 \iff 7x - 6x > 9 - 1 \iff \boxed{x > 8}.\]
Vậy tập nghiệm của bất phương trình là $S = \{x \in \mathbb{R} \mid x > 8\}$.
\end{loi_giai_cau}

\begin{de_cau}{Câu 3}{1,5 điểm}
\begin{enumerate}[label=\arabic*.]
\item Lớp 9B của một trường THCS có $40$ học sinh. Bảng thống kê sau cho biết số lượng học sinh của lớp 9B theo mức độ cận thị:
\begin{center}
\begin{tabular}{|c|c|c|c|}
\hline
Mức độ & Không cận thị & Cận thị nhẹ & Cận thị nặng \\ \hline
Số học sinh & $20$ & $12$ & $8$ \\ \hline
\end{tabular}
\end{center}
Lập bảng tần số tương đối cho bảng thống kê trên.
\item Tổ 1 của lớp 9A có 4 học sinh nam là An, Bình, Cường, Dũng và 5 học sinh nữ là Hoa, Lan, Mai, Ngọc, Phương. Giáo viên rút ngẫu nhiên một lá thăm để chọn một bạn tham gia hoạt động trải nghiệm.
Gọi $E$ là biến cố ``bạn được chọn là bạn nam'', $F$ là biến cố ``bạn được chọn là bạn nữ''.
\begin{enumerate}[label=\alph*)]
\item Có bao nhiêu kết quả có thể xảy ra khi chọn một bạn?
\item Tính xác suất của mỗi biến cố $E$ và $F$.
\end{enumerate}
\end{enumerate}
\end{de_cau}

\begin{loi_giai_cau}
\textbf{1.} Tổng số học sinh là $N = 40$.
Tần số tương đối của từng nhóm:
\begin{itemize}
\item Không cận thị: $f_1 = \dfrac{20}{40} \times 100\% = 50\%$.
\item Cận thị nhẹ: $f_2 = \dfrac{12}{40} \times 100\% = 30\%$.
\item Cận thị nặng: $f_3 = \dfrac{8}{40} \times 100\% = 20\%$.
\end{itemize}
Bảng tần số tương đối:
\begin{center}
\begin{tabular}{|c|c|c|c|c|}
\hline
Mức độ & Không cận thị & Cận thị nhẹ & Cận thị nặng & Tổng \\ \hline
Tần số tương đối & $50\%$ & $30\%$ & $20\%$ & $100\%$ \\ \hline
\end{tabular}
\end{center}

\textbf{2.}
\textbf{a)} Tổng số học sinh của tổ 1 là $4 + 5 = 9$ bạn.
Do đó có $\boxed{9}$ kết quả có thể xảy ra khi rút thăm ngẫu nhiên một bạn.

\textbf{b)}
\begin{itemize}
\item Biến cố $E$ có $4$ kết quả thuận lợi (An, Bình, Cường, Dũng), xác suất là: $P(E) = \boxed{\dfrac{4}{9}}$.
\item Biến cố $F$ có $5$ kết quả thuận lợi (Hoa, Lan, Mai, Ngọc, Phương), xác suất là: $P(F) = \boxed{\dfrac{5}{9}}$.
\end{itemize}
\end{loi_giai_cau}

\begin{de_cau}{Câu 4}{1,5 điểm}
\begin{enumerate}[label=\arabic*.]
\item Cho hàm số $y = ax^2$ ($a \neq 0$) có đồ thị Parabol đi qua điểm $A(2; 2)$. Tìm giá trị của $a$.
\item Cho phương trình: $x^2 - 4x + 1 = 0 \quad (1)$.
\begin{enumerate}[label=\alph*)]
\item Chứng minh rằng phương trình (1) luôn có hai nghiệm phân biệt $x_1, x_2$.
\item Không giải phương trình, hãy tính giá trị của biểu thức:
\[T = \dfrac{2}{x_2+1} + \dfrac{2}{x_1+1}.\]
\end{enumerate}
\end{enumerate}
\end{de_cau}

\begin{loi_giai_cau}
\textbf{1.} Đồ thị hàm số $y = ax^2$ đi qua điểm $A(2; 2)$ nên tọa độ điểm $A$ thỏa mãn công thức hàm số:
\[2 = a \cdot 2^2 \iff 4a = 2 \iff \boxed{a = \dfrac{1}{2}}.\]
Hàm số đã cho là $y = \dfrac{1}{2}x^2$.

\textbf{2.}
\textbf{a)} Phương trình $x^2 - 4x + 1 = 0$ có $a = 1, b' = -2, c = 1$.
Ta có: $\Delta' = (-2)^2 - 1 \cdot 1 = 4 - 1 = 3 > 0$.
Vì $\Delta' > 0$ nên phương trình (1) luôn có hai nghiệm phân biệt $x_1, x_2$.

\textbf{b)} Theo định lý Vi-ét:
\[\begin{cases} x_1 + x_2 = 4 \\ x_1 x_2 = 1 \end{cases}\]
Biến đổi biểu thức $T$:
\begin{align*}
T &= \frac{2}{x_2+1} + \frac{2}{x_1+1} = 2 \left(\frac{(x_1+1) + (x_2+1)}{(x_1+1)(x_2+1)}\right) \\
&= 2 \cdot \frac{(x_1+x_2) + 2}{x_1 x_2 + (x_1+x_2) + 1}.
\end{align*}
Thay $x_1+x_2 = 4$ và $x_1 x_2 = 1$ vào biểu thức:
\[T = 2 \cdot \frac{4 + 2}{1 + 4 + 1} = 2 \cdot \frac{6}{6} = \boxed{2}.\]
\end{loi_giai_cau}
